\chapter{Giới thiệu}
Chương này trình bày lý do lựa chọn đề tài, mục tiêu nghiên cứu, đối tượng và phạm vi thực hiện của khóa luận. Từ thực trạng dữ liệu y tế có tính phức tạp với các đặc tính phi tuyến, số chiều cao và mất cân bằng lớp, đề tài đặt ra vấn đề cần khảo sát hiệu quả của các mô hình học máy trong bài toán dự đoán trên hai bộ dữ liệu chăm sóc sức khỏe có đặc trưng riêng biệt. Phạm vi chính của khóa luận là đánh giá và so sánh hiệu suất của các mô hình học máy khác nhau, bao gồm cả các mô hình truyền thống và hiện đại, trong việc dự đoán kết quả lâm sàng trên hai bộ dữ liệu chăm sóc sức khỏe. 

\section{Lý do chọn đề tài}
Trong thập kỷ vừa qua, ngành y tế đang trải qua một cuộc chuyển đổi số sâu rộng. Sự phát triển của hồ sơ bệnh án điện tử (Electronic Health Records), công nghệ trình tự gene và các thiết bị theo dõi sức khỏe đã tạo ra một lượng dữ liệu khổng lồ mà các phương pháp phân tích truyền thống không còn đủ khả năng xử lý hiệu quả. Trong bối cảnh đó, các mô hình học máy nổi lên như một công cụ đầy tiềm năng để khai thác giá trị từ những nguồn dữ liệu này. Đơn cử như theo khảo sát của Bajwa và cộng sự (2021), trí tuệ nhân tạo và học máy đã được ứng dụng rộng rãi trong chẩn đoán bệnh, dự đoán nguy cơ bệnh lý, phân tích hình ảnh y tế và hỗ trợ ra quyết định lâm sàng.[1]\\

Tiềm năng ứng dụng của học máy trong y tế phụ thuộc rất lớn vào đặc tính của dữ liệu đầu vào. Khác với dữ liệu trong nhiều lĩnh vực khác, dữ liệu y tế thường tồn tại các đặc điểm phức tạp khiến các mô hình học máy thông thường không thể trực tiếp áp dụng mà không có bước tiền xử lý phù hợp. Trong phạm vi đề tài này, ba đặc tính được xác định là trực tiếp ảnh hưởng đến hai bộ dữ liệu được nghiên cứu, bao gồm: tính phi tuyến, số chiều cao (high-dimensional data) và mất cân bằng giữa các lớp.\\

Về số chiều cao, dữ liệu biểu hiện gen, loại dữ liệu phản ánh mức độ hoạt động của từng gen trong tế bào, là một trong những nguồn dữ liệu y tế phổ biến nhất trong nghiên cứu ung thư. Đặc điểm nổi bật của loại dữ liệu này là số lượng đặc trưng (gen) thường lên đến hàng chục nghìn, trong khi số lượng mẫu (bệnh nhân) chỉ ở mức vài trăm. Tình trạng đặc trưng lớn nhưng số lượng mẫu lại nhỏ này dẫn đến hiện tượng "lời nguyền đa chiều": khi không gian đặc trưng quá rộng so với lượng dữ liệu có sẵn, các mô hình học máy có xu hướng học thuộc dữ liệu huấn luyện thay vì nắm bắt quy luật thực sự, từ đó giảm đi khả năng tổng quát hóa trên dữ liệu mới của nó, hay còn gọi là hiện tượng học vẹt. Vì vậy, bước tiền xử lý , đặc biệt là lựa chọn đặc trưng đóng vai trò then chốt. Nghiên cứu của Bolon-Canedo và cộng sự đã chỉ ra rằng việc lựa chọn phương pháp chọn đặc trưng phù hợp có ảnh hưởng quyết định đến hiệu năng phân loại.[2] \\

Về tính phi tuyến, mối quan hệ giữa các đặc trưng sinh học và kết quả lâm sàng thường không thể biểu diễn bằng hàm tuyến tính đơn giản. Trong dữ liệu biểu hiện gen, các gen tương tác với nhau theo những cơ chế điều hòa phức tạp, tạo ra các mẫu phi tuyến trong không gian đặc trưng. Điều này đặt ra câu hỏi thực tiễn: liệu các mô hình phi tuyến như SVM với kernel RBF, Random Forest hay XGBoost có thực sự vượt trội so với mô hình tuyến tính như Logistic Regression trên dữ liệu y tế hay không, và ở mức độ nào?\\

Về mất cân bằng lớp, đây là vấn đề cố hữu của hầu hết bài toán phát hiện bệnh trong y tế: số ca bệnh (lớp thiểu số) luôn ít hơn nhiều so với ca không bệnh (lớp đa số). Sadeghi và cộng sự (2022) đã chỉ ra rằng mô hình học máy huấn luyện trên dữ liệu mất cân bằng thường đạt độ chính xác cao nhưng lại có độ bao phủ rất thấp trên lớp thiểu số, tức là bỏ sót phần lớn ca bệnh thực sự.[3] Đây là kết quả không thể chấp nhận được trong bối cảnh y tế vì việc bỏ sót chẩn đoán có thể ảnh hưởng trực tiếp đến tính mạng bệnh nhân.\\

Xuất phát từ những vấn đề trên, đề tài được thực hiện nhằm khảo sát một cách hệ thống hiệu quả của các mô hình học máy trên hai bộ dữ liệu y tế có đặc tính đại diện: bộ dữ liệu Breast Cancer Gene Expression với số chiều cao và tính phi tuyến, và bộ dữ liệu Diabetes Health Indicators với hiện tượng mất cân bằng lớp. Cụ thể, đề tài tập trung vào ba hướng khảo sát chính: (1) đánh giá tác động của feature selection đến hiệu năng phân loại trên dữ liệu số chiều cao, (2) đánh giá tác động của resampling đến khả năng phát hiện minority class trên dữ liệu mất cân bằng, và (3) trực quan hóa và so sánh hiệu năng mô hình trên cả hai bộ dữ liệu, nhằm làm rõ mức độ ảnh hưởng của từng đặc tính dữ liệu đến kết quả phân loại.
\section{Mục tiêu nghiên cứu}
Đề tài hướng đến bốn mục tiêu cụ thể sau:
\begin{itemize}
    \item Khảo sát và xử lý vấn đề số chiều cao trên bộ dữ liệu Breast Cancer Gene Expression (Biểu hiện gen của ung thư vú) thông qua hai nhóm phương pháp lựa chọn đặc trưng: .....
    \item Phân tích tính phi tuyến trên cả hai bộ dữ liệu thông qua trực quan hóa t-SNE và UMAP, kết hợp với phân tích đặc trưng quan trọng của các mô hình tree-based. Mức độ vượt trội của mô hình phi tuyến so với Logistic Regression sẽ được dùng làm bằng chứng định lượng cho tính phi tuyến của dữ liệu.
    \item Khảo sát ảnh hưởng của mất cân bằng lớp và so sánh hai kỹ thuật resampling — SMOTE và SMOTE+ENN — trên bộ dữ liệu Diabetes Health Indicators, đánh giá mức độ cải thiện F1-score và ROC-AUC so với baseline không resampling.
    \item So sánh hiệu năng của năm mô hình học máy trên cả hai pipeline: Logistic Regression, SVM (kernel RBF), Random Forest, XGBoost và Multi-layer Perceptron (MLP). Mục tiêu là xác định mô hình nào ổn định nhất trong bối cảnh dữ liệu y tế có đặc tính phức tạp.    
\end{itemize}

\section{Đối tượng và phạm vi nghiên cứu}

Đề tài sử dụng hai bộ dữ liệu healthcare được lấy từ các kho dữ liệu công khai, mỗi bộ đại diện cho một nhóm thách thức đặc thù của dữ liệu y tế.

\subsection{Breast Cancer Gene Expression Dataset (GSE45827 - CuMiDa)}

Bộ dữ liệu GSE45827 được thu thập từ cơ sở dữ liệu CuMiDa (Curated Microarray Database), một kho dữ liệu microarray đã được làm sạch và tuyển chọn từ GEO nhằm phục vụ cho các nghiên cứu và đánh giá mô hình học máy trong lĩnh vực ung thư. Bộ dữ liệu gồm 151 mẫu bệnh nhân, trong đó mỗi mẫu được mô tả bởi 54.676 đặc trưng gen. Các mẫu được chia thành nhiều nhóm tương ứng với các phân nhóm ung thư vú khác nhau, bao gồm TNBC, HER2, Luminal A, Luminal B và mô bình thường. Một đặc điểm quan trọng của bộ dữ liệu là số lượng đặc trưng rất lớn so với số lượng mẫu. Cụ thể, trung bình mỗi mẫu có khoảng 362 đặc trưng gene trên một quan sát, tạo nên bài toán có số lượng đặc trưng lớn, số lượng mẫu nhỏ. Trong các bài toán như vậy, mô hình học máy thường gặp khó khăn do dữ liệu có độ chiều cao, dễ xảy ra hiện tượng overfitting và làm giảm khả năng tổng quát hóa trên dữ liệu mới. Vì vậy, việc áp dụng các kỹ thuật lựa chọn đặc trưng và giảm chiều dữ liệu là cần thiết trước khi xây dựng mô hình dự đoán.

\subsection{Diabetes Health Indicators Dataset (BRFSS 2015 - CDC)}

Bộ dữ liệu được xây dựng từ cuộc khảo sát BRFSS 2015 (Behavioral Risk Factor Surveillance System) do CDC (Centers for Disease Control and Prevention) thực hiện hằng năm nhằm thu thập thông tin về hành vi sức khỏe và các bệnh mãn tính của người dân Mỹ. Bộ dữ liệu gồm 253.680 mẫu khảo sát với 21 đặc trưng đầu vào, bao gồm các chỉ số lâm sàng và hành vi như BMI, huyết áp, cholesterol, mức độ hoạt động thể chất, thói quen hút thuốc và tiền sử bệnh tim mạch. Nhãn lớp phân thành 2 nhóm: không mắc tiểu đường và tiểu đường, với mức độ mất cân bằng rõ rệt giữa hai nhóm này (tỉ lệ xấp xỉ 86:14). Theo thống kê của IDF (2021), có khoảng 537 triệu người mắc tiểu đường trên toàn cầu và con số này được dự báo tăng lên 783 triệu vào năm 2045,[5] cho thấy tầm quan trọng của các mô hình dự đoán sớm cho căn bệnh này.

\subsection{Phạm vi nghiên cứu}

Để đảm bảo tính khả thi và chiều sâu của thực nghiệm, đề tài được giới hạn trong các phạm vi sau:

\begin{itemize}
    \item Chỉ sử dụng các phương pháp Supervised Learning cho bài toán phân loại. Các phương pháp \textit{Unsupervised Learning} như clustering hay anomaly detection không phải là đối tượng đánh giá chính của nghiên cứu.

    \item Không nghiên cứu chuyên sâu về Deep Learning. Mô hình MLP được sử dụng như một mạng nơ-ron nông nhằm làm đường cơ sở để so sánh với các mô hình học máy truyền thống, thay vì tối ưu hóa kiến trúc mạng nơ-ron.

    \item Không triển khai hệ thống trong môi trường bệnh viện thực tế. Toàn bộ thực nghiệm được thực hiện trong môi trường nghiên cứu bằng ngôn ngữ Python với các thư viện chính gồm \texttt{scikit-learn}, \texttt{XGBoost}, \texttt{imbalanced-learn}, \texttt{matplotlib} và \texttt{seaborn}.

    \item Kỹ thuật lựa chọn đặc trưng chỉ được áp dụng cho bộ dữ liệu Breast Cancer, trong khi các phương pháp tái lấy mẫu được áp dụng cho bộ dữ liệu Diabetes nhằm xử lý hiện tượng mất cân bằng lớp. Việc khảo sát tính phi tuyến được thực hiện trên cả hai pipeline nghiên cứu.
\end{itemize}

\subsection{Nội dung thực hiện}

Khóa luận được tổ chức thành hai pipeline thực nghiệm độc lập, mỗi pipeline tập trung giải quyết một nhóm vấn đề đặc thù của dữ liệu y tế. Nội dung thực hiện bao gồm các bước chính sau:

\begin{itemize}
    \item \textbf{Tiền xử lý dữ liệu:} xử lý giá trị thiếu, chuẩn hóa dữ liệu bằng StandardScaler và mã hóa nhãn. Tất cả các bước tiền xử lý được huấn luyện trên tập train và áp dụng sang tập test.

    \item \textbf{Trực quan hóa và phân tích tính phi tuyến:} sử dụng các kỹ thuật giảm chiều t-SNE và UMAP để chiếu dữ liệu xuống không gian hai chiều, từ đó quan sát khả năng phân tách giữa các lớp và đánh giá đặc tính phi tuyến của dữ liệu. Bước này được áp dụng cho cả hai pipeline.

    \item \textbf{Lựa chọn đặc trưng (Pipeline Breast Cancer):} so sánh các phương pháp thuộc nhóm filter gồm \texttt{Variance Threshold} và \texttt{SelectKBest} với phương pháp wrapper \texttt{RFECV}. Hiệu năng mô hình được đánh giá trên ba phiên bản dữ liệu gồm dữ liệu gốc, dữ liệu sau filter và dữ liệu sau wrapper.

    \item \textbf{Xử lý mất cân bằng lớp (Pipeline Diabetes):} so sánh hiệu quả của các phương pháp \texttt{SMOTE} và \texttt{SMOTE+ENN} với trường hợp không sử dụng tái lấy mẫu trên năm mô hình phân loại.

    \item \textbf{Huấn luyện và tối ưu mô hình:} năm mô hình gồm Logistic Regression, SVM RBF, Random Forest, XGBoost và MLP được huấn luyện bằng phương pháp \textit{Stratified K-Fold Cross-Validation} ($k=5$) kết hợp với \textit{GridSearchCV} để tối ưu siêu tham số.

    \item \textbf{Đánh giá và thảo luận kết quả:} so sánh hiệu năng các mô hình dựa trên các độ đo Accuracy, F1-score, ROC-AUC và Confusion Matrix. Đồng thời phân tích tầm quan trọng của đặc trưng (\textit{feature importance}) để đối chiếu với kết quả trực quan hóa và đặc điểm của dữ liệu.
\end{itemize}

\subsection{Bố cục khóa luận}

Khóa luận được tổ chức thành năm chương như sau:

\begin{itemize}
    \item \textbf{Chương 1 -- Giới thiệu:} trình bày bối cảnh nghiên cứu, lý do lựa chọn đề tài, mục tiêu nghiên cứu và phạm vi thực hiện.

    \item \textbf{Chương 2 -- Cơ sở lý thuyết:} trình bày các kiến thức nền tảng liên quan đến dữ liệu y tế, các kỹ thuật trực quan hóa dữ liệu, lựa chọn đặc trưng, các mô hình phân loại, kỹ thuật xử lý mất cân bằng lớp và các độ đo đánh giá mô hình.

    \item \textbf{Chương 3 -- Phương pháp nghiên cứu:} mô tả chi tiết quy trình thực nghiệm của từng pipeline, thiết lập thực nghiệm, quy trình xử lý dữ liệu và các câu hỏi nghiên cứu được đặt ra.

    \item \textbf{Chương 4 -- Kết quả thực nghiệm:} trình bày và phân tích kết quả đạt được trên từng pipeline, đồng thời thực hiện so sánh tổng hợp giữa các pipeline nghiên cứu.

    \item \textbf{Chương 5 -- Kết luận và hướng phát triển:} tổng kết các kết quả chính của khóa luận, nêu các hạn chế còn tồn tại và đề xuất những hướng nghiên cứu tiếp theo.
\end{itemize}

